\PassOptionsToPackage{table,xcdraw}{xcolor}
\documentclass[11pt, a4paper]{article}

\usepackage[T1]{fontenc}
\usepackage[utf8]{inputenc}
\usepackage{tgpagella}
\usepackage[spanish, es-noshorthands]{babel}
\usepackage{amsmath, amssymb}
\usepackage{booktabs}
\usepackage{tabularx}
\usepackage[most]{tcolorbox}
\usepackage[margin=2.5cm]{geometry}
\usepackage{fancyhdr}

\pagestyle{fancy}
\fancyhf{}
\setlength{\headheight}{16pt}
\lhead{\textbf{UCB Sede Tarija} | Documento Docente}
\chead{Notas Técnicas: Lab 4 Físico}
\rhead{Circuitos Electrónicos II}
\renewcommand{\headrulewidth}{0.4pt}
\definecolor{ucbblue}{RGB}{0, 51, 102}

\begin{document}

\begin{center}
    \Large\textbf{\textcolor{ucbblue}{Guía Técnica Docente: Lab 4 Físico (Dinámica)}}[cite: 13]
\end{center}

\begin{tcolorbox}[colback=white, colframe=black, title=\textbf{Diferenciación Crítica de Conceptos[cite: 13]}]
    \begin{itemize}
        \item \textbf{Predicción:} Estimación teórica calculada por el estudiante[cite: 13].
        \item \textbf{Valor Típico Datasheet:} Un parámetro nominal del fabricante, \textbf{no una medida física garantizada}[cite: 13].
        \item \textbf{Resultado Físico:} Lo que realmente mide el osciloscopio bajo tolerancia, temperatura e instrumentación[cite: 13].
    \end{itemize}
\end{tcolorbox}

\section*{1. Predicciones Teóricas (Referencias Típicas de Fabricante)[cite: 13]}
\textbf{Comparación Principal Sugerida:} Familia 741 (LM741 o UA741CN) vs TL074[cite: 13]. Si hay restricciones de inventario: LM358N vs TL074[cite: 13]. ¡Debe verificarse el fabricante exacto del stock antes de asignar el datasheet![cite: 13]

\begin{table}[h!]
    \centering
    \footnotesize
    \renewcommand{\arraystretch}{1.3}
    \begin{tabularx}{\textwidth}{|l|c|c|X|X|}
    \hline
    \textbf{Dispositivo / Familia} & \textbf{$GBW_{\mathrm{typ}}$} & \textbf{$SR_{\mathrm{typ}}$} & \textbf{Exp A: $f_{BW}$ ($A_N=10$)} & \textbf{Exp B: Umbral SR ($V_p=4\,\text{V}$)} \\ \hline
    Familia 741 (LM/UA)[cite: 13] & $\approx 1\,\text{MHz}$ & $\approx 0.5\,\text{V}/\mu\text{s}$ & $\approx 100\,\text{kHz}$ & $\approx 19.9\,\text{kHz}$ \\ \hline
    TI LM358 (Estándar)[cite: 13] & $\approx 0.7\,\text{MHz}$ & $\approx 0.3\,\text{V}/\mu\text{s}$ & $\approx 70\,\text{kHz}$ & $\approx 11.9\,\text{kHz}$ \\ \hline
    TI TL074 (Quad FET)[cite: 13] & $\approx 3\,\text{MHz}$ & $\approx 13\,\text{V}/\mu\text{s}$ & $\approx 300\,\text{kHz}$ & $\approx 517\,\text{kHz}$ \\ \hline
    \end{tabularx}
\end{table}

\section*{2. Secuencia de Diagnóstico (Troubleshooting)[cite: 13]}
Si un grupo obtiene salidas inesperadas, aplique este orden de descarte \textbf{antes} de evaluar GBW o SR[cite: 13]:
\begin{enumerate}
    \setlength{\itemsep}{0pt}
    \item ¿Rieles correctos ($\pm10\,\text{V}$) y medidos con referencia a $0\,\text{V}$ (punto medio si usan baterías)?[cite: 13]
    \item ¿Pinout exacto (DIP-8 vs DIP-14)? ¿Intercambio inadvertido de entradas $+$ y $-$?[cite: 13]
    \item ¿Ganancia DC correcta a $1\,\text{kHz}$?[cite: 13]
    \item ¿Amplitud de entrada leída del osciloscopio o confiada a ciegas en el display del generador?[cite: 13]
    \item ¿Existe recorte (clipping) a baja frecuencia por amplitud excesiva?[cite: 13]
    \item ¿Limitación de ancho de banda del propio generador o punta de prueba atenuada inadvertidamente?[cite: 13]
\end{enumerate}

\section*{3. Preguntas de Verificación Oral[cite: 13]}
\begin{itemize}
    \item \textit{¿Por qué el Exp A usa señal pequeña y el Exp B reduce la ganancia?} (Para aislar los efectos de GBW y SR, respectivamente)[cite: 13].
    \item \textit{¿Qué evidencia visual sugiere Slew Rate en lugar de Clipping?} (Triangularización con pendiente constante vs Picos chatos con riel saturado)[cite: 13].
    \item \textit{¿Qué valor de datasheet usaste: típico, mínimo o máximo?} (Evaluar consciencia de tolerancia estadística del silicio)[cite: 13].
    \item \textit{¿Por qué LTSpice y el osciloscopio podrían no coincidir?} (El macromodelo es una aproximación y puede no incluir limitaciones de gran señal precisas)[cite: 13].
\end{itemize}

\end{document}
